\documentclass[11pt]{article}
\usepackage[a4paper,margin=25mm]{geometry}
\usepackage{amsmath,amssymb,amsthm,booktabs,array}
\usepackage[T1]{fontenc}
\usepackage{lmodern}
\usepackage[hidelinks]{hyperref}
\newtheorem{theorem}{Theorem}
\newtheorem{lemma}[theorem]{Lemma}
\newtheorem{corollary}[theorem]{Corollary}
\newtheorem{proposition}[theorem]{Proposition}
\newcommand{\ECN}{E_{\mathrm{CN}}}
\newcommand{\CN}{\mathrm{CN}}
\title{Proofs of Nath--Fasfous--Das--Shang's commuting and surface conjectures, and a sharp form of Tak--Dutta--Nath's clique theorem}
\author{Haobo Ma\\\small ChronoAI Pte Ltd, Singapore
\and Wenlin Zhang\\\small National University of Singapore\\\small ChronoAI Pte Ltd, Singapore}
\date{17 November 2026}
\begin{document}
\maketitle
\begin{abstract}
For Nath--Fasfous--Das--Shang's Conjectures~1--3
(\emph{Symmetry} \textbf{13} (2021), 1651, Section~4, p.~11), we prove
the surface and commuting-graph statements and give an alternative proof
with exact equality cases for Tak--Dutta--Nath's prior clique result
(arXiv:2605.22160v1 (2026), Theorem~5.4, pp.~14--15).
Every planar or toroidal graph is non-CN-hyperenergetic.
For every finite nonabelian group, its noncentral commuting graph satisfies
$\ECN\le n(n-z-2)<2(n-1)(n-2)$, where $n$ counts noncentral elements
and $z$ counts central elements; the same bound holds for every finite
noncommutative ring, without an identity assumption.
The group argument combines the published estimate $\ECN\le4m$ with a
proper-centralizer count. A neighborhood-clique decomposition proves the
surface statement. For a union of cliques, equality with the complete-graph
benchmark occurs exactly for one clique or two isolated vertices.
\end{abstract}

\section{The question and the spectral bounds}
Nath, Fasfous, Das and Shang~\cite{NFDS} conjectured that planar or
toroidal graphs, noncentral commuting graphs of finite nonabelian groups,
and disjoint unions of cliques are not CN-hyperenergetic.
We prove the first two assertions and extend the group bound to all finite
noncommutative rings. The clique assertion has a prior proof by Tak, Dutta
and Nath~\cite[Theorem~5.4]{TDN}; Section~\ref{sec:cliques} gives another
proof and identifies the exceptions to the strict inequality displayed there.
Jannat, Nath and Das~\cite[Remark~4.7, equation~(6), p.~15]{JND} supply
the general estimate $\ECN(H)\le4|E(H)|$ used for the group result.

For a finite simple graph $H$, let $A$ be its adjacency matrix, $D$ its
diagonal degree matrix, and $N_H(v)$ the neighbors of $v$. Following
Alwardi, Soner and Gutman~\cite[Definitions~1.5--1.7]{ASG}, put
\[
 \CN(H)_{uv}=\begin{cases}|N_H(u)\cap N_H(v)|,&u\ne v,\\0,&u=v,\end{cases}
 \qquad \ECN(H)=\sum_i|\lambda_i(\CN(H))|.
\]
All eigenvalues are counted with multiplicity. Counting two-step walks gives
$\CN(H)=A^2-D$~\cite[Proposition~2.7]{ASG}.
For $n\ge2$, $\CN(K_n)=(n-2)(J_n-I_n)$, where $J_n$ is the all-ones
matrix and $I_n$ is the identity. Thus the comparison value is
\[
 B(n)=\ECN(K_n)=2(n-1)(n-2) \quad(n\ge1).
\]
The value at $n=1$ is zero. A graph of positive order $n$ is
CN-hyperenergetic when $\ECN(H)>B(n)$.

\begin{lemma}\label{lem:bounds}
If $H$ has $m$ edges and $r$ vertices of positive degree, then
\[
 \ECN(H)\le4m,\qquad
 \ECN(H)\le2\sum_v\max\{d_H(v)-1,0\}=4m-2r.
\]
\end{lemma}
\begin{proof}
For the first inequality, set $P=A^2$ and $Q=D$, both positive
semidefinite. In an orthonormal eigenbasis $u_i$ of $P-Q$, put
$p_i=u_i^{\mathsf T}Pu_i\ge0$ and $q_i=u_i^{\mathsf T}Qu_i\ge0$.
Then
\[
 \sum_i|\lambda_i(P-Q)|=\sum_i|p_i-q_i|
 \le\sum_i(p_i+q_i)=\operatorname{tr}P+\operatorname{tr}Q=4m.
\]
For the second inequality, let $B_v$ have an entry~1 at distinct
vertices of $N_H(v)$ and zero elsewhere. Entrywise counting gives
$\CN(H)=\sum_v B_v$. If $d=d_H(v)\ge1$, the nonzero block of $B_v$
is $J_d-I_d$, whose eigenvalues are $d-1$ and $-1$ with multiplicity
$d-1$. Its trace norm, the sum of its singular values, is $2(d-1)$;
for $d=0$ it is zero. The trace norm of a real symmetric matrix equals
the sum of its absolute eigenvalues. Its triangle inequality
\cite[Lemma~4.2]{JND} now gives
the result, since $\sum_v d_H(v)=2m$.
\end{proof}

\section{All finite nonabelian groups and noncommutative rings}
For a group $G$, write $Z(G)$ for its center and $C_G(x)$ for the
centralizer of $x$. The graph $\Gamma_c(G)$ has the individual elements
of $G\setminus Z(G)$ as vertices; distinct $x,y$ are adjacent when
$xy=yx$. Define $\Gamma_c(R)$ in the same way for a ring $R$.
Rings here are associative, with an abelian additive group; a
multiplicative identity is not required.

\begin{theorem}\label{thm:groups}
Let $G$ be a finite nonabelian group, $z=|Z(G)|$, and
$n=|G\setminus Z(G)|$. Then $z\ge1$, $n\ge z+2\ge3$, and
\begin{equation}\label{eq:main}
 \ECN(\Gamma_c(G))\le n(n-z-2)<B(n).
\end{equation}
In particular, Conjecture~2 of~\cite[Section~4, p.~11]{NFDS} holds.
\end{theorem}
\begin{proof}
Put $N=|G|=n+z$. For noncentral $x$, the subgroup $C_G(x)$ is proper,
so its index is at least two and $|C_G(x)|\le N/2$. The disjoint union
\[
 C_G(x)=Z(G)\,\dot\cup\,\{x\}\,\dot\cup\,N_{\Gamma_c(G)}(x)
\]
therefore gives
\begin{equation}\label{eq:degree}
 d(x)=|C_G(x)|-z-1,\qquad 2d(x)\le n-z-2.
\end{equation}
There is a noncentral element, and its degree is nonnegative; hence
$n\ge z+2$. Also $z\ge1$ because the identity is central.
By Lemma~\ref{lem:bounds},
$\ECN(\Gamma_c(G))\le2\sum_xd(x)\le n(n-z-2)$.
Finally,
\[
 B(n)-n(n-z-2)=(n-2)^2+nz>0.
\]
\end{proof}

\begin{corollary}\label{cor:rings}
If $R$ is any finite noncommutative ring, $z=|Z(R)|$, and
$n=|R\setminus Z(R)|$, then $n\ge z+2\ge3$ and
\[
 \ECN(\Gamma_c(R))\le n(n-z-2)<B(n).
\]
\end{corollary}
\begin{proof}
The additive centralizer $C_R(x)=\{y\in R:xy=yx\}$ is a proper
additive subgroup for noncentral $x$. Closure follows from
$x(y-w)=xy-xw=(y-w)x$. Its order is at most $|R|/2$.
The disjoint decomposition into $Z(R)$, $\{x\}$ and the neighbors of
$x$ is unchanged, and $0\in Z(R)$. The proof of
Theorem~\ref{thm:groups} applies word for word.
\end{proof}
Fasfous and Nath~\cite[Conjecture~3.12, p.~75]{FN} posed the
non-hyperenergetic conclusion for finite rings whose noncentral
centralizers are commutative (CC-rings). Tak, Dutta and
Nath~\cite[Corollary~5.5, p.~15]{TDN} proved that case.
Corollary~\ref{cor:rings} removes the CC assumption; the extension to
all finite noncommutative rings is not posed in either statement.

\begin{corollary}\label{cor:prime}
For either $S=G$ or $S=R$ as above, let $p$ be the least prime divisor
of $|S|$. Then
\[
 \ECN(\Gamma_c(S))\le
 2\sum_{x\notin Z(S)}\max\{|C_S(x)|-z-2,0\},\qquad
 \ECN(\Gamma_c(S))\le2n\left(\frac{n+z}{p}-z-1\right).
\]
\end{corollary}
\begin{proof}
The first bound is the second inequality of Lemma~\ref{lem:bounds}
and the degree identity~\eqref{eq:degree}.
The index of each proper centralizer subgroup divides $|S|$, so it
is at least $p$. Replace $|S|/2$ by $|S|/p$ in the degree argument.
\end{proof}

\section{Planar and toroidal graphs}
We first dispose of the small orders without enumerating all graphs.

\begin{lemma}\label{lem:small}
Every simple graph on at most seven vertices and of positive order is
non-CN-hyperenergetic.
\end{lemma}
\begin{proof}
Isolated vertices do not affect the energy, and $B(n)$ is nondecreasing
for positive integers. It therefore suffices to consider graphs with no
isolated vertices; the cases of at most two vertices have energy zero.
For orders $n=3,4,5,6,7$, respectively, Lemma~\ref{lem:bounds} gives
$\ECN\le4m-2n\le B(n)$ whenever
\[
 m\le2,\ 5,\ 8,\ 13,\ 18,
\]
respectively. The remaining graphs, up to isomorphism, are $K_n$
($3\le n\le7$), $K_n-e$ ($5\le n\le7$), and $K_7$ with two edges
removed, either disjoint or incident at a common endpoint.

Each remaining common-neighborhood matrix has exactly one positive
eigenvalue. Here are the details. For $K_n$ this follows from
$\CN(K_n)=(n-2)(J_n-I_n)$. For the other graphs partition the vertices
as indicated below. The subspace constant on each part is invariant;
the matrix acting on these constants is the quotient matrix $Q$.
The remaining eigenvalues and the quotients are
\[
\begin{array}{c|c|c}
\text{removed edges; part sizes}&Q&\text{remaining eigenvalues}\\ \hline
e;\ 2,n-2 &
\begin{pmatrix}n-2&(n-2)(n-3)\\2(n-3)&(n-2)(n-3)\end{pmatrix}
&[-(n-2)]^{n-2}\\[8pt]
2K_2;\ 4,3&\begin{pmatrix}11&12\\16&10\end{pmatrix}&[-5]^4,[-1]^1\\[8pt]
K_{1,2};\ 1,2,4&\begin{pmatrix}0&8&12\\4&4&16\\3&8&15\end{pmatrix}
&[-4]^1,[-5]^3
\end{array}
\]
For $K_n-e$, the first part consists of the missing edge's endpoints;
for $2K_2$ it contains all four endpoints; for $K_{1,2}$ the parts are
the common endpoint, the two other endpoints, and the four remaining
vertices. Within a part, zero-sum vectors give the displayed eigenvalues;
in the $2K_2$ case, the contrast of the two endpoint-pair sums gives $-1$.
These vectors and the constant-on-part subspace span the whole space.
The first two quotients have negative determinants
$(n-2)(n-3)(4-n)$ and $-82$. Their eigenvalues therefore have opposite
signs. The last quotient has characteristic polynomial
\[
 t^3-19t^2-136t-144.
\]
Its roots are real because a quotient of a symmetric matrix for such a
partition is similar to a symmetric matrix. Their product is $144>0$,
and their pairwise-product sum is $-136<0$, so exactly two are negative.
Thus there is exactly one positive eigenvalue in every case.

For any zero-trace symmetric matrix with one positive eigenvalue $\rho$,
its energy is $2\rho$. Each row sum of $\CN(H)$ is at most
$(n-1)(n-2)$. The spectral radius satisfies the same bound, so
$2\rho\le B(n)$, as required.
\end{proof}

\begin{theorem}\label{thm:surfaces}
If $H$ is a finite simple planar or toroidal graph of positive order $n$,
then $\ECN(H)\le B(n)$. Thus Conjecture~1 of
\cite[Section~4, p.~11]{NFDS} holds.
\end{theorem}
\begin{proof}
Delete isolated vertices, leaving $r$ vertices and $m$ edges. If $r\le7$,
apply Lemma~\ref{lem:small}; the empty remainder has energy zero.
If $r\ge8$, Euler's edge bound for a simple graph embeddable in the
torus gives $m\le3r$; this also holds for planar graphs. To see this,
connect components by additional edges and choose a minimum-genus
cellular embedding, of genus $g\le1$. If the augmented graph has $m'$
edges and $f$ faces, then $r-m'+f=2-2g$ and $3f\le2m'$, so
$m\le m'\le3r-6+6g\le3r$.
By Lemma~\ref{lem:bounds}, $\ECN\le4m-2r\le10r$, while
\[
 B(r)-10r=2(r-4)^2-28\ge4>0.
\]
Adding back isolated vertices can only increase the benchmark.
\end{proof}

\section{The clique theorem and its equality cases}\label{sec:cliques}
The non-strict conclusion of the following proposition is Conjecture~3
of~\cite[Section~4, p.~11]{NFDS}, proved by Tak, Dutta and
Nath~\cite[Theorem~5.4, pp.~14--15]{TDN}.

\begin{proposition}\label{prop:cliques}
If $H=K_{s_1}\sqcup\cdots\sqcup K_{s_k}$, with $k\ge1$, $s_i\ge1$
and $n=\sum_i s_i$, then
\[
 \ECN(H)=2\sum_i(s_i-1)(s_i-2)\le B(n).
\]
Equality holds exactly for $k=1$, or $k=2$ and $s_1=s_2=1$.
\end{proposition}
\begin{proof}
The matrix is block diagonal across components, so the energies add.
Put $f(s)=2(s-1)(s-2)$ for $s\ge1$. For positive integers $a,b$,
\[
 f(a+b)-f(a)-f(b)=4(ab-1)\ge0,
\]
with equality exactly when $a=b=1$. Merge components successively.
If $k\ge3$, a second merger involving the first merged part is strict,
since that part has at least two vertices. This gives both directions
of the equality characterization.
\end{proof}
The unrestricted strict inequality displayed in~\cite[Theorem~5.4]{TDN}
therefore needs these exceptions: a single clique, for example $K_3$,
and $2K_1$ give equality. The source conjecture itself asks only for
non-hyperenergeticity and is true as stated.

\section{Coverage and mathematical evidence}
\begin{center}\small
\begin{tabular}{@{}p{0.30\textwidth}p{0.64\textwidth}@{}}
\toprule Source item & Coverage\\\midrule
\cite{NFDS}, Conjecture 1 & Proved: Theorem~\ref{thm:surfaces}.\\[3pt]
\cite{NFDS}, Conjecture 2 & Proved: Theorem~\ref{thm:groups}, with a center-dependent bound.\\[3pt]
\cite{NFDS}, Conjecture 3 & Reproved; prior proof Tak--Dutta--Nath (2026).
Proposition~\ref{prop:cliques} gives all equality cases.\\[3pt]
\cite{TDN}, Theorem 5.4, strict form & Refuted-and-corrected by
Proposition~\ref{prop:cliques}; equality for $K_3$ and $2K_1$.\\[3pt]
\cite{FN}, Conjecture 3.12 & Reproved; prior proof Tak--Dutta--Nath (2026),
Corollary 5.5. Corollary~\ref{cor:rings} extends it to all finite
noncommutative rings (an extension not posed in that source).\\[3pt]
\cite{NFDS}, Theorems 7--17 and Corollaries 4--6 & Their non-hyperenergetic
conclusions follow from Theorem~\ref{thm:groups}; exact energy formulas
are not rederived.\\[3pt]
\cite{ASG}, Open problems 1--2 & Not covered in general: existence of
CN-hyperenergetic graphs and effects of edge deletion. Lemma~\ref{lem:small}
excludes the former through order seven.\\[3pt]
\cite{TDN}, Problems 5.1--5.2 & Not covered: integrality and hyperenergeticity
for the distinct minimum-second-neighborhood matrix.\\
\bottomrule\end{tabular}\end{center}

\paragraph{Formal statements and written proofs.}
The declaration \texttt{NoncentralBound.result} in
\href{https://github.com/the-omega-institute/trureturing/blob/b3bd002bc644799939f299fe6c6c26195e3bb998/D5/S3/Combinatorics/Graph/CommonNeighborhood/NoncentralBound.lean}{\texttt{NoncentralBound.lean}}
states $\ECN(\Gamma_c(G))\le B(n)$ for every finite nonabelian group.
Its helpers \texttt{cn\_energy\_le\_twice\_degree\_sum} and
\texttt{twice\_degree\_bound} express the spectral and centralizer estimates;
the center-dependent estimate occurs in the proof of \texttt{result}.
The companion file
\href{https://github.com/the-omega-institute/trureturing/blob/b3bd002bc644799939f299fe6c6c26195e3bb998/D5/S3/Combinatorics/Graph/SubgroupGenerating/CyclicExtensionRefutation.lean}{\texttt{CyclicExtensionRefutation.lean}}
defines \texttt{commonNeighbourMatrix} with a zero diagonal,
\texttt{graphSpectrum} with the full eigenvalue multiset, and
\texttt{graphEnergy} with zero centering for CN-energy.
The strict group comparison, ring and prime-divisor bounds,
surface theorem, and clique equality cases have the written proofs above;
they are distinct from the stated scope of \texttt{result}.

\paragraph{Exact finite checks.}
The accompanying standard-library programs check the matrix identity,
clique equality cases, finite group and ring centralizers, and the quotient
matrices in Lemma~\ref{lem:small} using integers and rational numbers.
The quotient characteristic polynomials and invariant subspaces are checked
directly. The universal conclusions follow from the written proofs.
The supporting JSON files specify the finite ranges and exact outputs.

\begin{thebibliography}{9}
\bibitem{NFDS} R.~K.~Nath, W.~N.~T.~Fasfous, K.~C.~Das and Y.~Shang,
Common Neighborhood Energy of Commuting Graphs of Finite Groups,
\emph{Symmetry} \textbf{13} (2021), 1651.
\href{https://doi.org/10.3390/sym13091651}{doi:10.3390/sym13091651}.
\bibitem{JND} Firdous Ee Jannat, R.~K.~Nath and K.~C.~Das,
Common neighborhood energies and their relations with Zagreb index,
\href{https://arxiv.org/abs/2402.15416v1}{arXiv:2402.15416v1} (2024).
\bibitem{ASG} A.~Alwardi, N.~D.~Soner and I.~Gutman,
On the common-neighborhood energy of a graph,
\emph{Bull. Acad. Serbe Sci. Arts, Sci. Math.} \textbf{36} (2011), 49--59.
\href{https://www.maths.tcd.ie/EMIS/journals/BSANU/36/r4.pdf}{Full text}.
\bibitem{TDN} P.~Tak, J.~Dutta and R.~K.~Nath,
Minimum second neighborhood degree energy of commuting graphs of finite rings,
\href{https://arxiv.org/abs/2605.22160v1}{arXiv:2605.22160v1} (2026).
\bibitem{FN} W.~N.~T.~Fasfous and R.~K.~Nath,
Common neighborhood spectrum and energy of commuting graphs of finite rings,
\emph{Palestine J. Math.} \textbf{13}(1) (2024), 66--76.
\href{https://pjm.ppu.edu/sites/default/files/papers/PJM_13\%281\%29_2024_66_to_76.pdf}{Full text}.
\end{thebibliography}
\medskip
\noindent\small This work is licensed under the Creative Commons Attribution 4.0
International license (CC BY 4.0).
\end{document}
